\section{Problema}
\label{subsec:problema}

Un problema (computacional, evidentemente) es un enunciado que especifica qué se quiere obtener, sin decir cómo.
\begin{definicion}{Problema (computacional)}{problema}
Especificación de una entrada y de una salida esperada con base en la entrada.
\end{definicion}

E.g., el \concept{problema de ordenamiento} (\textit{sorting problem}, muy famoso):
\begin{itemize}
  \item \textbf{Entrada}: secuencia de \(n\) números \(a_1, a_2,\ldots,a_n\).
  \item \textbf{Salida}: permutación \(a_1', a_2',\ldots,a_n'\) tal que \(a_1' \leq a_2' \leq \ldots \leq a_n'\)
\end{itemize}
Recuérdese que cada permutación de una secuencia es un posible orden de sus elementos. En este caso, la salida esperada es cualquier permutación en la que los números estén en orden no decreciente.

% Hmm no sé si me gusta este apartado
En ocasiones los problemas están redactados de forma que la entrada y la salida esperada están ocultas en el enunciado. En estas notas ese nunca será el caso, pues no se pretende ejercitar la habilidad de interpretar un problema, solo la de resolverlo.

Una \concept{instancia} de un problema es una entrada específica, e.g., \(1, 4, 2, 3, 2\) es una instancia del problema de ordenamiento, en cuyo caso una \concept{solución} es \(1, 2, 2, 3, 4\).

\subsection{Tipos de problemas}

Los problemas se caracterizan en tipos con base en la estructura de la salida esperada. Un mismo problema de la vida real se puede modelar de varias formas, que pueden corresponder a diferentes tipos. Algunos tipos son:
\begin{itemize}
  \item \concept{Enumeración}: la salida son todas las soluciones válidas.
  \item \concept{Conteo}: la salida es el número de soluciones válidas.
  \item \concept{Optimización}: la salida es la mejor solución válida de acuerdo a algún criterio.
  \item \concept{Búsqueda}: la salida es alguna solución válida para el problema, si existe.
  \item \concept{Decisión}\label{problema_decision}: la salida es sí o no.
\end{itemize}
Esas categorías no son mutuamente excluyentes, frecuentemente son distintas perspectivas del mismo problema. E.g., el problema de ordenamiento se puede abordar bajo cualquiera de las perspectivas anteriores:
\begin{itemize}
  \item Enumeración: dado un arreglo, listar todas las permutaciones ordenadas de sus elementos.
  \item Conteo: dado un arreglo, determinar cuántas permutaciones distintas de sus elementos están ordenadas.
  \item Optimización: dado un arreglo y un conjunto de operaciones permitidas, encontrar una secuencia de operaciones que lo ordene usando el menor número de operaciones posible.
  \item Búsqueda: dado un arreglo, encontrar alguna permutación de sus elementos que esté ordenada.
  \item Decisión: dado un arreglo, ¿está ordenado?
\end{itemize}
Se pueden formular aún más problemas relacionados a ordenar elementos que encajen en cualquiera de esas categorías. 
